% Generated draft: requires mathematical and source review.
\documentclass[11pt]{article}
\usepackage[margin=25mm]{geometry}
\usepackage{fontspec}
\setmainfont{Latin Modern Roman}
\usepackage{amsmath,amssymb,mathtools,unicode-math}
\setmathfont{Latin Modern Math}
\usepackage{enumitem,needspace,xurl,hyperref}
\setlength{\parindent}{0pt}
\setlength{\parskip}{5pt}
\setlength{\emergencystretch}{3em}
\newcommand{\N}{\mathbb{N}}
\newcommand{\Z}{\mathbb{Z}}
\newcommand{\Q}{\mathbb{Q}}
\newcommand{\R}{\mathbb{R}}
\newcommand{\C}{\mathbb{C}}
\newcommand{\F}{\mathbb{F}}
\newcommand{\A}{\mathbb{A}}
\newcommand{\PP}{\mathbb{P}}
\newcommand{\E}{\mathbb{E}}
\newcommand{\eps}{\varepsilon}
\newcommand{\dd}{\,\mathrm{d}}
\begin{document}
\section*{UnsolvedMath ID: 30002377}
\textbf{Source problem code:} catalogue entry
\subsection{Definitions and mathematical statement}
Let \(\mathbb T=\R/\Z\). A frequency \(\omega\in\mathbb T\) is Diophantine if there exist constants \(\kappa>0\) and \(\tau>0\) such that
\[
\operatorname{dist}(k\omega,\Z)\geq \kappa |k|^{-\tau}
\qquad\text{for every }k\in\Z\setminus\{0\},
\]
where the distance is computed using any real representative of \(\omega\).

The supplied baseline assumptions are \(f\in C^2(\mathbb T,\R)\) and the existence of distinct points \(c_1,c_2\in\mathbb T\) such that
\[
f'(x)=0\iff x\in\{c_1,c_2\},\qquad f''(c_1)\ne0,\qquad f''(c_2)\ne0.
\]
The frequency \(\omega\) is Diophantine, and the coupling \(\lambda\) is sufficiently large in the sense intended by the source.

For the transfer matrix \(A_E(x)\) of the referenced quasiperiodic Schr\"odinger cocycle, its induced projective map has the form
\[
\Phi_E(x,[v])=\bigl(x+\omega,[A_E(x)v]\bigr),
\qquad (x,[v])\in\mathbb T\times\mathbb{RP}^{1},
\]
where \([v]\) denotes the line spanned by \(v\in\R^2\setminus\{0\}\). Identifying \(\mathbb{RP}^{1}\) with a circle gives the stated torus domain \(\mathbb T^2\). The supplied record does not specify the transfer-matrix normalization or the projective coordinate.

\emph{Source question.} Describe the dynamics of \(\Phi_E\) for every \(E\in\R\), extending the setting above to more general assumptions on \(f\).

This is an unresolved question schema rather than a fully specified quantified conjecture: the record does not identify which assumptions on \(f\) are to be weakened, which replacement hypotheses are intended, or which dynamical properties constitute the requested description. It also does not specify the coupling threshold, its parameter dependence, or whether one threshold must work simultaneously for every energy. The baseline assumptions therefore cannot be silently treated as the requested generalized class.
\subsection{Short English statement}
The question asks for a description, at every real energy, of the projective dynamics of a quasiperiodic Schr\"odinger cocycle at sufficiently large coupling, under assumptions on the sampling function broader than the supplied \(C^2\), two-nondegenerate-critical-point setting. The intended broader class and the required dynamical conclusions are unspecified, so a unique precise conjecture cannot be recovered from this record alone.
\subsection{Sources}
\begin{itemize}
\item Local UnsolvedMath registry: supplied complete record for permanent ID \(30002377\), titled \emph{Dynamics of Large-Coupling Quasiperiodic Schr\"odinger Cocycles}; exact source problem code \texttt{catalogue entry}, in the Oberwolfach Reports Open Problems collection. No individual proposer is supplied; the recorded proposed year is 2013.
\item Supplied report citation: \emph{Dynamische Systeme (2014)}, Oberwolfach Reports, DOI \texttt{10.4171/owr/2013/34}, pp.~1975--2033. \url{https://doi.org/10.4171/owr/2013/34}. The background describes Oberwolfach Reports as documenting mathematical workshops and their open-problem discussions.
\item Supplied literature source: \emph{Oberwolfach Report 34/2013}. \url{https://doi.org/10.4171/owr/2013/34}.
\item Literature assessment dated exactly \texttt{2026-08-21}: supplied dated curation, not independently verified here. It labels the problem open, states that the report leaves the dynamics under more general assumptions on \(f\) open, and says that available results cover analytic, Gevrey, and selected differentiable classes rather than the requested generality. The supplied background expressly calls for independent verification of time-sensitive status and cited claims.
\item The registry labels the statement \texttt{corrected\_verified} and reports that its mathematical content was checked against the official report, with notation, context, and question form normalized without changing the source-level claim. This is supplied verification metadata, not verification performed for this draft.
\end{itemize}
\end{document}
